\documentclass[a4paper,11pt]{article}
\usepackage[utf8]{inputenc}
\usepackage[T1]{fontenc}
\usepackage[spanish]{babel}
\usepackage{geometry}
\usepackage{array}
\usepackage{colortbl}
\usepackage{xcolor}
\usepackage{booktabs}
\usepackage{enumitem}
\usepackage{fancyhdr}
\usepackage{tikz}

\geometry{left=2cm, right=2cm, top=2.5cm, bottom=2cm}

\definecolor{violeta}{HTML}{7B2D8E}
\definecolor{azulg}{HTML}{3498DB}
\definecolor{rojog}{HTML}{E74C3C}
\definecolor{gris}{HTML}{F2F2F2}
\definecolor{verde}{HTML}{27AE60}
\definecolor{naranja}{HTML}{E67E22}

\pagestyle{fancy}
\fancyhf{}
\fancyhead[L]{\small\textbf{Nuevo Compañeros 1 — U03: La Familia}}
\fancyhead[R]{\small\textbf{Píldora formativa 3.10}}
\fancyfoot[C]{\thepage}

\setlength{\parindent}{0pt}

\begin{document}

\begin{center}
{\LARGE\textbf{PÍLDORA FORMATIVA 3.10}}\\[0.3cm]
{\Large Robot Dreams: del dato cultural a la reflexión intercultural}\\[0.2cm]
{\normalsize Sección Cultura (p.42) — Proyectable — 3 diapositivas}\\[0.2cm]
{\small Versión enriquecida secuencial para Fases CU2--CU5}
\end{center}

\vspace{0.5cm}

\noindent\fbox{\parbox{\dimexpr\linewidth-2\fboxsep-2\fboxrule}{%
\textbf{Contenido para el profesor}\\[0.2cm]
El texto de la página 42 es un texto divulgativo que presenta una introducción al cómic como forma artística y una presentación de \textit{Robot Dreams}, película de animación española dirigida por Pablo Berger (2023). La película es muda (sin diálogos), narra la amistad entre un perro solitario y un robot en Manhattan, fue nominada al Óscar a la Mejor Película de Animación (2024) y ganó el Goya (2024).\\[0.2cm]
El texto funciona como \textbf{contenedor lingüístico}: contiene 7 verbos en presente regular ya formalizados, vocabulario de la unidad y solo 6 palabras nuevas (profesiones del cine). Al hacer visible este reciclaje con colores, el texto se vuelve transparente y el foco pasa al contenido cultural.
}}

\vspace{0.8cm}

%% ============================================================
%% DIAPOSITIVA 1
%% ============================================================

\noindent\textcolor{violeta}{\rule{\linewidth}{2.5pt}}

\section*{DIAPOSITIVA 1 — EL TEXTO CULTURAL COMO CONTENEDOR}

{\small Fase CU2 (Lectura global + texto mapeado) — Técnica: Texto cromático con mapeo de reciclaje lingüístico}\\
{\small\textit{Principio:} El texto cultural no es un texto nuevo que el estudiante debe descifrar palabra a palabra. Contiene gramática y vocabulario ya formalizados. Al hacerlos visibles con colores, el texto se vuelve transparente y el foco pasa al contenido cultural. El estudiante gana confianza: ya sé mucho de lo que dice este texto.}

\vspace{0.5cm}

\subsection*{En pantalla}

\begin{center}
{\Large\textbf{¿QUÉ SABÉIS YA DE ESTE TEXTO?}}
\end{center}

\vspace{0.3cm}

El texto cultural de la página 42 se presenta con tres categorías de color:

\vspace{0.3cm}

\noindent\textcolor{verde}{\rule{0.5cm}{0.5cm}}\hspace{0.3cm}\textbf{\textcolor{verde}{VERDE — Verbos en presente regular ya formalizados:}}

\vspace{0.2cm}

\begin{center}
\renewcommand{\arraystretch}{1.5}
\begin{tabular}{|>{\centering\arraybackslash}p{3.5cm}|>{\centering\arraybackslash}p{3.5cm}|>{\centering\arraybackslash}p{3.5cm}|}
\hline
\rowcolor{verde!10}
\textbf{Verbo en el texto} & \textbf{Infinitivo} & \textbf{Conjugación} \\
\hline
\textcolor{verde}{\textbf{nace}} & nacer & -er \\
\hline
\textcolor{verde}{\textbf{necesita}} & necesitar & -ar \\
\hline
\textcolor{verde}{\textbf{intervienen}} & intervenir & -ir \\
\hline
\textcolor{verde}{\textbf{vive}} & vivir & -ir \\
\hline
\textcolor{verde}{\textbf{fabrica}} & fabricar & -ar \\
\hline
\textcolor{verde}{\textbf{se separan}} & separarse & -ar \\
\hline
\textcolor{verde}{\textbf{interesan}} & interesar & -ar \\
\hline
\end{tabular}
\end{center}

\vspace{0.3cm}

\noindent\textcolor{azulg}{\rule{0.5cm}{0.5cm}}\hspace{0.3cm}\textbf{\textcolor{azulg}{AZUL — Vocabulario y estructuras de la unidad:}}

\begin{itemize}[leftmargin=2cm, itemsep=2pt]
\item \textcolor{azulg}{\textbf{historia}} (vocabulario general)
\item \textcolor{azulg}{\textbf{amigo inseparable}} (vocabulario de relaciones)
\item \textcolor{azulg}{\textbf{de todas las edades}} (expresión de descripción)
\end{itemize}

\vspace{0.3cm}

\noindent\textcolor{rojog}{\rule{0.5cm}{0.5cm}}\hspace{0.3cm}\textbf{\textcolor{rojog}{ROJO — Vocabulario temático nuevo (profesiones del cine):}}

\begin{itemize}[leftmargin=2cm, itemsep=2pt]
\item \textcolor{rojog}{\textbf{dibujantes}}, \textcolor{rojog}{\textbf{guionistas}}, \textcolor{rojog}{\textbf{iluminadores}}, \textcolor{rojog}{\textbf{maquilladores}}, \textcolor{rojog}{\textbf{actores}}, \textcolor{rojog}{\textbf{decoradores}}
\end{itemize}

\vspace{0.3cm}

\begin{center}
\fbox{\parbox{13cm}{\centering\large
\textcolor{verde}{7 verbos en presente regular} + \textcolor{azulg}{vocabulario de la unidad} + \textcolor{rojog}{6 profesiones nuevas}\\[0.2cm]
{\normalsize El texto está lleno de cosas que ya sabéis.}
}}
\end{center}

\vspace{0.5cm}

\subsection*{Instrucciones para el profesor}

\begin{enumerate}[leftmargin=1.5cm, itemsep=4pt]
\item Después de la primera lectura global (ejercicio 2, CU2), proyecte o escriba en la pizarra las tres categorías de color.
\item Pida a los estudiantes que busquen en el texto: \textit{¿Cuántos verbos en presente regular encontráis?} Recoja respuestas a mano alzada. Subraye cada verbo en verde.
\item Señale el vocabulario ya conocido en azul: \textit{historia, amigo}. Pregunte: \textit{¿Dónde habéis visto estas palabras antes?}
\item Señale en rojo las profesiones del cine: estas son las palabras nuevas. Solo 6. Y ya las habéis trabajado en el ejercicio 1.
\item Cierre el mapeo: \textit{Este texto tiene 7 verbos que ya conocéis, vocabulario de la unidad y solo 6 palabras nuevas. Ya sabéis mucho de lo que dice. Ahora vamos a leer QUÉ nos cuenta sobre el cine.}
\end{enumerate}

\newpage

%% ============================================================
%% DIAPOSITIVA 2
%% ============================================================

\noindent\textcolor{violeta}{\rule{\linewidth}{2.5pt}}

\section*{DIAPOSITIVA 2 — ROBOT DREAMS: MÁS ALLÁ DEL DATO}

{\small Fase CU4 (Explotación del contenido cultural específico) — Técnica: Reflexión intercultural guiada}\\
{\small\textit{Principio:} El libro presenta datos sobre \textit{Robot Dreams} (director, argumento, premios). La píldora transforma esos datos en experiencia intercultural. La película es el vehículo: desde ella, el estudiante reflexiona sobre cómo la animación conecta culturas, cómo una película puede ser universal sin usar palabras, y cómo el cine español va más allá de los estereotipos.}

\vspace{0.5cm}

\subsection*{En pantalla}

\begin{center}
{\Large\textbf{ROBOT DREAMS — ¿POR QUÉ ES ESPECIAL?}}
\end{center}

\vspace{0.5cm}

\noindent\colorbox{violeta!10}{\parbox{\dimexpr\linewidth-2\fboxsep}{%
\textbf{\textcolor{violeta}{1. UNA PELÍCULA SIN PALABRAS}}\\[0.2cm]
\textit{Robot Dreams} es una película muda. No tiene diálogos.\\[0.1cm]
¿Por qué un director elige no usar palabras?\\
¿Qué puede decir una película sin hablar?
}}

\vspace{0.4cm}

\noindent\colorbox{azulg!10}{\parbox{\dimexpr\linewidth-2\fboxsep}{%
\textbf{\textcolor{azulg}{2. TEMAS UNIVERSALES}}\\[0.2cm]
La amistad, la soledad, el amor.\\[0.1cm]
¿Estos temas son iguales en tu cultura?\\
¿Hay películas de animación en tu país que hablen de estos temas?
}}

\vspace{0.4cm}

\noindent\colorbox{naranja!10}{\parbox{\dimexpr\linewidth-2\fboxsep}{%
\textbf{\textcolor{naranja}{3. LA PARADOJA}}\\[0.2cm]
Los críticos dicen que es \textit{la película más humana de la temporada}, pero entre sus personajes no hay un solo humano.\\[0.1cm]
¿Cómo puede ser la más humana sin tener humanos?
}}

\vspace{0.5cm}

\begin{center}
\fbox{\parbox{12cm}{\centering
\textbf{Recuerda:} diferente no es mejor ni peor. Diferente es diferente.
}}
\end{center}

\vspace{0.5cm}

\subsection*{Instrucciones para el profesor}

\begin{enumerate}[leftmargin=1.5cm, itemsep=4pt]
\item Antes de proyectar, haga la transición desde CU3: \textit{Ya sabéis qué dice el texto. Ahora vamos a pensar sobre lo que dice.}
\item Lea la primera pregunta. Deje un momento de silencio para pensar. Recoja 2--3 respuestas. Acepte respuestas en L1 si el nivel lingüístico no permite la expresión en español.
\item Lea la segunda pregunta. En parejas, 2 minutos para comentar. Recoja comparaciones: \textit{¿Qué películas de animación de vuestro país conocéis?}
\item Lea la tercera pregunta (la paradoja). Esta es la más rica. Guíe: los personajes son un perro y un robot, pero hablan de cosas muy humanas. A veces la animación dice más que una película con actores reales.
\item Matice generalizaciones si surgen: si un estudiante dice \textit{Los españoles hacen películas raras}, reformule: \textit{Algunas películas españolas son diferentes a las de tu país. Pero no todas. ¿Conoces otras películas españolas?}
\end{enumerate}

\newpage

%% ============================================================
%% DIAPOSITIVA 3
%% ============================================================

\noindent\textcolor{violeta}{\rule{\linewidth}{2.5pt}}

\section*{DIAPOSITIVA 3 — COMPARAR SIN JERARQUIZAR}

{\small Fases CU4--CU5 (Explotación cultural + Producción) — Técnica: Aplicación de la estrategia intercultural}\\
{\small\textit{Principio:} El estudiante aplica la estrategia de la tarjeta (Caja 3) para comparar el cine de su cultura con el cine hispano. La comparación se hace sin jerarquizar: no hay cine mejor ni peor, hay cine diferente.}

\vspace{0.5cm}

\subsection*{En pantalla}

\begin{center}
{\Large\textbf{COMPARAR SIN JERARQUIZAR}}
\end{center}

\vspace{0.5cm}

\begin{center}
\renewcommand{\arraystretch}{1.6}
\begin{tabular}{|>{\raggedright\arraybackslash}p{3.5cm}|>{\raggedright\arraybackslash}p{5.5cm}|>{\raggedright\arraybackslash}p{5.5cm}|}
\hline
\rowcolor{violeta!15}
 & \textbf{\textcolor{violeta}{En la cultura hispana}} & \textbf{\textcolor{violeta}{En mi cultura}} \\
\hline
\textbf{Películas de animación} & \textit{Robot Dreams} (España), \textit{Metegol} (Argentina), \textit{Anina} (Uruguay)\ldots & ¿Cuáles conoces de tu país? \\
\hline
\textbf{Cómics} & \textit{Mafalda} (Argentina), \textit{Mortadelo y Filemón} (España)\ldots & ¿Cuáles conoces de tu país? \\
\hline
\textbf{Temas} & Amistad, soledad, familia, humor\ldots & ¿Son los mismos temas? ¿Diferentes? \\
\hline
\end{tabular}
\end{center}

\vspace{0.5cm}

\noindent\colorbox{verde!10}{\parbox{\dimexpr\linewidth-2\fboxsep}{%
\textbf{\textcolor{verde}{Expresiones para comparar sin jerarquizar:}}\\[0.3cm]
\begin{itemize}[leftmargin=1cm, itemsep=3pt]
\item En España hay películas de animación como \textit{Robot Dreams}. En mi país hay\ldots
\item Las dos culturas tienen películas de animación, pero son diferentes porque\ldots
\item Me gusta \textit{Robot Dreams} porque\ldots{} / En mi país me gusta\ldots{} porque\ldots
\item No todas las películas españolas / de mi país son iguales.
\end{itemize}
}}

\vspace{0.5cm}

\begin{center}
\fbox{\parbox{12cm}{\centering
\textbf{Recuerda los 4 pasos de la tarjeta:}\\[0.2cm]
{\large Observa $\rightarrow$ Compara $\rightarrow$ Matiza $\rightarrow$ Valora}
}}
\end{center}

\vspace{0.5cm}

\subsection*{Instrucciones para el profesor}

\begin{enumerate}[leftmargin=1.5cm, itemsep=4pt]
\item Reparta la tarjeta de estrategia intercultural (Caja 3) si no lo ha hecho antes.
\item Lea los 4 pasos de la tarjeta. Modele con un ejemplo personal: \textit{A mí me gustan las películas de animación japonesas. Muchas películas japonesas hablan de la naturaleza. Robot Dreams habla de la amistad. Son diferentes, pero las dos me gustan.}
\item En parejas: cada estudiante completa la tabla con ejemplos de su cultura (2--3 minutos).
\item Puesta en común: cada pareja comparte una similitud y una diferencia. El profesor recoge en la pizarra.
\item Cierre con matización: si alguien ha generalizado (\textit{En España no hay buenas películas de animación} / \textit{Las películas de mi país son mejores}), reformule usando las expresiones de la tarjeta: \textit{Algunas películas españolas son diferentes a las de tu país. ¿Pero son peores? No, son diferentes.}
\end{enumerate}

\end{document}
